\section{Odd kernels and a third-order balancing-Pell packet}
\label{sec:pell-odd-kernel-third-order}

Continue with
$(1+\sqrt2)^n=A_n+B_n\sqrt2$ and an odd prime index $\ell$.  Suppose, for
pressure testing, that both $A_\ell$ and $B_\ell$ are squarefull.  A finite
squarefull factorization has the canonical form
\[
 N=D^3C^2,\qquad
 D=\prod_{v_p(N)\text{ odd}}p.
\]
Indeed an even exponent is $2+2c$ and an odd exponent is $3+2c$.  The
inherited perfect-power classifications force each exponent vector in a
hypothetical packet to have gcd one, so both odd kernels are nontrivial:
\begin{equation}\label{eq:pell-odd-kernel-normal-form}
 A_\ell=D_A^3X^2,\qquad B_\ell=D_B^3Y^2,\qquad
 D_A,D_B>1.
\end{equation}
Every kernel prime consequently occurs at depth at least three in its
assigned channel.  The negative Pell equation becomes
\[
                   D_A^6X^4-2D_B^6Y^4=-1.
\]

Let $s_\ell=(2/\ell)$.  The rank congruences, sign product, and residue orbit
from the preceding Pell sections give
\begin{align}
 D_A&\equiv1\pmod{2\ell},&
 D_B&\equiv s_\ell\pmod{2\ell},
 \label{eq:pell-kernel-rank}\\
 X^2&\equiv1\pmod{2\ell},&Y^2&\equiv1\pmod{2\ell}.
 \label{eq:pell-core-square}
\end{align}
Because $\ell$ is prime and the square cores are odd, each of $X,Y$ is
congruent to one of $\pm1$ modulo $2\ell$.  The complete mod-eight data are
\[
\begin{array}{c|rrrr}
\ell\bmod8&1&3&5&7\\ \hline
D_A\bmod8&1&7&1&7\\
D_B\bmod8&1&5&5&1.
\end{array}
\]
Writing $d_A=(D_A-1)/(2\ell)$ and
$d_B=(D_B-s_\ell)/(2\ell)$ further gives
\[
\begin{array}{c|rrrr}
\ell\bmod8&1&3&5&7\\ \hline
d_A\bmod4&0&1&0&1\\
d_B\bmod4&0&1&3&0.
\end{array}
\]

The two kernels also form a three-edge Jacobi-character triangle with the
index:
\[
 \left(\frac{D_A}{D_B}\right)=s_\ell,\qquad
 \left(\frac{\ell}{D_A}\right)=\left(\frac{-1}{\ell}\right),
 \qquad
 \left(\frac{\ell}{D_B}\right)=
       \left(\frac{s_\ell}{\ell}\right).
\]
Thus the three signs, in the order displayed, are
\[
\begin{array}{c|rrrr}
\ell\bmod8&1&3&5&7\\ \hline
(D_A/D_B)&+1&-1&-1&+1\\
(\ell/D_A)&+1&-1&+1&-1\\
(\ell/D_B)&+1&-1&+1&+1.
\end{array}
\]
In the $3\pmod8$ class all three aggregate edges are negative.  This does
not assert that every individual cross pair is negative; an actual index-
eleven factorization already refutes that stronger statement.

One additional quotient digit can be retained.  For a finite list $T$, let
$E_j(T)$ be its $j$th elementary symmetric sum.  The exact expansion
\[
 \prod_i(1+xt_i)\equiv
 1+xE_1(T)+x^2E_2(T)+x^3E_3(T)\pmod{x^4}
\]
follows by induction.  Repeat each $A$-channel quotient according to its
valuation to form $T_A$, and repeat each signed $B$-channel quotient to form
$T_B$.  Write $(K_A,C_A,H_A)=(E_1,E_2,E_3)(T_A)$ and similarly for $B$.
For
\[
 a=\frac{A_\ell-1}{2\ell},\qquad
 b=\frac{s_\ell B_\ell-1}{2\ell},
\]
cancellation of $2\ell$ gives
\[
 a\equiv K_A+2\ell C_A+4\ell^2H_A,\qquad
 b\equiv K_B+2\ell C_B+4\ell^2H_B\pmod{8\ell^3}.
\]
Substitution into the exact identity
$a-2b+\ell(a^2-2b^2)=0$ yields the necessary third-order ledger
\begin{align*}
0\equiv{}&K_A-2K_B+\ell(K_A^2-2K_B^2)+2\ell(C_A-2C_B)\\
 &+4\ell^2(H_A-2H_B+K_AC_A-2K_BC_B)\\
 &+4\ell^3(C_A^2-2C_B^2)\pmod{8\ell^3}.
\end{align*}
Its reduction modulo $4\ell^2$ recovers the preceding second-order ledger.
The new congruence is a necessary condition and is not asserted to be a
contradiction.

A deterministic two-direction search examined all 668 odd prime indices
$3\le\ell\le5000$ and all primes $q\le2{,}000{,}000$ in the necessary
classes $q\equiv\pm1\pmod{2\ell}$.  It certified 481 indices with an actual
exponent-one divisor of $A_\ell$ or $B_\ell$; 187 bounded-search rows remain
open.  An independent verifier recomputed every witness modulo $q^2$, checked
primality, and, for every odd prime $\ell\le191$, verified a simple divisor
from the exact coordinates.  The largest such witness,
\[
        13558774610046711780701\parallel B_{59},
\]
is supported by a complete Pocklington certificate.  Hence
\[
 \boxed{A_\ell B_\ell\text{ is not squarefull for every odd prime }
        \ell\le191.}
\]
This is a finite theorem.  It neither proves eventual simplicity nor treats
an unresolved row as a squarefull packet.

The Lean companion checks the cubic-square factor identity, residue
transport, Jacobi reciprocity steps, arbitrary-list third-order expansion,
quotient cancellation, and the coupled ledger.  The remaining positive gate
is an actual opposite-channel depth-three exclusion satisfying all kernel,
residue, character, and quotient constraints.  No full-premise packet
counterexample was found, so this route remains active.
